% 05_state_prep.tex
\section{State preparation on unbounded domains}
\label{sec:stateprep}

The no-go of Sec.~\ref{sec:limits} closes the door on chart change
within block encoding.  We now open a different one.  Pairing the
stereographic chart with a cotangent grid and a diagonal
Fourier-type signal operator prepares a quantum state encoding a
function $f\colon\R\to\C$ on the whole real line, at a cost set only
by the shape of $f$ after pullback to the circle.  The signal
operator is a walk with prescribed eigenphases, not a block
encoding of a Hermitian operator, so Proposition~\ref{prop:nogo}
does not apply.

\subsection{The cotangent grid}
\label{subsec:grid}

We prepare an $N=2^n$-dimensional state
\begin{equation}
	\ket{\Psi_f}=\frac{1}{\mathcal{N}_f}\sum_{j=0}^{N-1}f(y_j)\ket{j},
	\qquad \mathcal{N}_f=\Bigl(\textstyle\sum_j|f(y_j)|^2\Bigr)^{1/2},
	\label{eq:target-state}
\end{equation}
on the \emph{cotangent grid}
\begin{equation}
	y_j=L\cot\theta_j,\qquad \theta_j=\frac{\pi(2j+1)}{2N},
	\qquad j=0,\dots,N-1,
	\label{eq:cot-grid}
\end{equation}
the image under $y=L\cot\theta$ of a uniform angular grid.  The
half-integer shift keeps every point interior to $(0,\pi)$.  The
map parameter $L>0$ sets the density of points near $y=0$; the grid
clusters at the origin and thins in the tails, the reverse of the
Chebyshev grid on $[-1,1]$.  Its extent is $|y|\lesssim 2LN/\pi$,
which grows exponentially with $n$.  The cotangent grid thus ties
the represented range to the qubit count $n$; it does not
\emph{eliminate} truncation, but replaces a user-chosen radius by a
range fixed by the resource $n$, beyond which the unrepresented
tail mass is exponentially small in $n$ for the targets below.

\subsection{The Boyd class}
\label{subsec:boyd}

Given $f\colon\R\to\C$, its stereographic pullback is
$g(\theta)=f(L\cot\theta)$.  Because $\cot$ has period $\pi$, $g$ is
$\pi$-periodic, so its Fourier series contains only the harmonics
$\cos 2k\theta$ and $\sin 2k\theta$.  The rational Chebyshev
functions of Eq.~\eqref{eq:rat-cheb} satisfy
$\TB_k(y_j;L)=\cos k\theta_j$ and $\SB_k(y_j;L)=\sin k\theta_j$, so
a degree-$d$ rational Chebyshev approximation of $f$ on $\R$ is a
degree-$d$ trigonometric approximation of $g$ on the circle.

\begin{definition}[Boyd class~\cite{Boyd1990}]
	\label{def:boyd}
	$f$ belongs to $\mathcal{B}(L)$ if it is free of singularities
	on the real axis and its pullback $g$ has Fourier coefficients
	obeying $|a_k|,|b_k|\le C\exp(-q k^\sigma)$ for some $C,q>0$ and
	$\sigma\in(0,1]$.
\end{definition}

By Boyd's theorem~\cite{Boyd1990}, $f\in\mathcal{B}(L)$ with
\emph{geometric} rate ($\sigma=1$) if and only if $g$ extends
analytically to a neighborhood of the real $\theta$-axis, which
requires $f$ analytic near the real $y$-axis \emph{and} analytic at
infinity, i.e., possessing a convergent expansion in $1/y$.  The
degree for $L^\infty$-error $\delta$ is then
$d_\delta=\mathcal{O}(\log(1/\delta))$.  When $f$ has an essential
singularity at infinity the pullback has an essential singularity
at the endpoint $\theta=0$, the analyticity strip collapses, and
the rate is \emph{subgeometric}, $\sigma<1$, with
$d_\delta=\mathcal{O}((\log(1/\delta))^{1/\sigma})$.  Two targets
of interest are subgeometric with $\sigma=1/2$: the Gaussian, whose
essential singularity at infinity is standard, and the hyperbolic
secant.  For $\mathrm{sech}(L\cot\theta)$ the poles of $\mathrm{sech}$
pull back to points $\theta_m$ with $|\im\theta_m|\to 0$ as
$m\to\infty$, accumulating on the real axis, so no strip exists and
the coefficients decay as $\exp(-c\sqrt{k})$
(Appendix~\ref{app:sech}).  The hyperbolic secant is therefore
\emph{not} in the geometric class, contrary to a natural reading of
its exponential decay; exponential decay on $\R$ is necessary but
not sufficient for geometric convergence, which needs analyticity
at infinity.

\subsection{Filling fraction}
\label{subsec:filling}

The success probability of the preparation is governed by the
filling fraction
\begin{equation}
	\mathcal{F}_g^{[N]}=\frac{1}{\sqrt{N}}\cdot
	\frac{\bigl(\sum_{j=0}^{N-1}|g(\theta_j)|^2\bigr)^{1/2}}{|g|_{\max}},
	\qquad |g|_{\max}=\max_{(0,\pi)}|g|.
	\label{eq:filling}
\end{equation}

\begin{lemma}[Continuous limit]
	\label{lem:filling}
	If $g\in C^2(0,\pi)$ with $|g|_{\max}<\infty$ attained at an
	interior point, then
	\begin{equation}
		\mathcal{F}_g^{[N]}=\mathcal{F}_g^{[\infty]}+\mathcal{O}(N^{-2}),
		\qquad
		\mathcal{F}_g^{[\infty]}=\frac{\|g\|_{L^2(0,\pi)}}
		{\sqrt{\pi}\,|g|_{\max}}.
		\label{eq:filling-limit}
	\end{equation}
\end{lemma}

The sum in Eq.~\eqref{eq:filling} is the midpoint rule for
$(1/\pi)\int_0^\pi|g|^2$ on the grid $\theta_j$, whose error is
$\mathcal{O}(N^{-2})$ for $g\in C^2$; the square root and the
$\mathcal{O}(N^{-2})$ agreement of the discrete and continuous
maxima preserve the rate (Appendix~\ref{app:filling}).  The
continuous limit depends only on the shape of $g$ and on no
truncation parameter.  For the targets below, $g$ is $C^\infty$
with bounded derivatives, so the discrete and continuous filling
fractions agree to $1+\mathcal{O}(2^{-n})$.

\subsection{Algorithm and main result}
\label{subsec:algorithm}

The algorithm has three stages.  \emph{Classical preprocessing}
computes the pullback $g$, normalizes to $\hat g=g/|g|_{\max}$ with
$|\hat g|\le 1$, and approximates $\hat g$ by a degree-$d_\delta$
trigonometric polynomial $\tilde g$ to $L^\infty$-error $\delta$;
because $g$ is $\pi$-periodic, $\tilde g$ is built from even
harmonics on the full circle, so $|\tilde g|\le 1+\delta$ holds on
all of $\T$ and is restored to $|\tilde g|\le 1$ by rescaling with
$1/(1+\delta)$.  The QSP phases follow from any standard
solver~\cite{DongMengWhaleyEtAl2021,BerntsonSunderhauf2024,NiSarkarYingLin2025}
at cost negligible against the circuit for the targets here.
\emph{The Fourier-QET circuit} uses the diagonal signal operator
\begin{equation}
	e^{iA}=\sum_{j=0}^{N-1}e^{i\theta_j}\ket{j}\!\bra{j},\qquad
	\theta_j=\frac{\pi(2j+1)}{2N},
	\label{eq:qet-signal}
\end{equation}
whose phase $\theta_j=\pi/(2N)+\sum_{k}j_k\pi 2^k/N$ decomposes over
the bits of $j$, so $e^{iA}$ costs $\mathcal{O}(n)$ elementary
gates: $n$ single-qubit $R_z$ rotations, one on each qubit $k$ with
angle $\pi 2^k/N$, and one global phase.  With
$d_\delta$ applications of $e^{\pm iA}$ and $\mathcal{O}(d_\delta)$
ancilla rotations, the circuit block-encodes the diagonal
$D_{\tilde g}=\sum_j\tilde g(\theta_j)\ket{j}\!\bra{j}$.
\emph{Amplitude amplification} boosts the flagged block to unity in
$\mathcal{O}(1/\mathcal{F}_{\tilde g}^{[N]})$ rounds with one extra
ancilla.

\begin{theorem}[Stereographic state preparation]
	\label{thm:stateprep}
	Let $f\in\mathcal{B}(L)$ and let a degree-$d_\delta$
	trigonometric polynomial approximate $\hat g$ to $L^\infty$-error
	$\delta=\tfrac12\varepsilon\cdot
	\min(\mathcal{F}_g^{[N]},\mathcal{F}_{\tilde g}^{[N]})$ on the
	circle.  The algorithm prepares a state $\varepsilon$-close in
	trace distance to $\ket{\Psi_f}$ using
	\begin{equation}
		\mathcal{O}\!\left(\frac{n\,d_\delta}{\mathcal{F}_{\tilde g}^{[N]}}\right)
		\text{ gates}
		\label{eq:stateprep-cost}
	\end{equation}
	and at most two ancilla qubits, one for the signal and one for
	amplitude amplification.  For the geometric class
	$d_\delta=\mathcal{O}(\log(1/\varepsilon))$ and the cost is
	$\mathcal{O}(n\log(1/\varepsilon)/\mathcal{F}_g^{[\infty]})$;
	for the subgeometric class
	$d_\delta=\mathcal{O}((\log(1/\varepsilon))^{1/\sigma})$.
\end{theorem}

The proof (Appendix~\ref{app:stateprep}) applies the QSP
completeness theorem on the full circle, evaluates the resulting
diagonal on $\ket{+}^{\otimes n}$, and converts the $L^\infty$
approximation error to trace distance through the filling fraction;
the choice $\delta=\tfrac12\varepsilon\min(\mathcal{F}_g^{[N]},
\mathcal{F}_{\tilde g}^{[N]})$ controls both filling fractions
without assuming an order between them.  A degree-$d_\delta$
trigonometric polynomial has $2d_\delta+1\le N$ degrees of freedom
on the $N$ grid points, so no aliasing enters while
$d_\delta=\mathcal{O}((\log(1/\varepsilon))^{1/\sigma})$ stays below
$N/2$, which holds at every qubit count of interest.

\subsection{Examples}
\label{subsec:examples}

For the Lorentzian $f(y)=1/(1+y^2)$ with $L=1$, the pullback is
\begin{equation}
	g(\theta)=\frac{1}{1+\cot^2\theta}=\sin^2\theta
	=\tfrac12(1-\cos 2\theta),
\end{equation}
a trigonometric polynomial of degree two with no approximation
error, $\delta=0$.  The filling fraction is
$\mathcal{F}_g^{[\infty]}=\sqrt{3/8}\approx 0.612$, one round of
amplitude amplification suffices, and the cost is $\mathcal{O}(n)$.
More generally $1/(1+y^2)^s$ for integer $s$ pulls back to
$\sin^{2s}\theta$, a trigonometric polynomial of degree $2s$, so the
whole family has cost independent of the represented range.

For the hyperbolic secant, $\sigma=1/2$ (Sec.~\ref{subsec:boyd}),
so $d_\delta=\mathcal{O}(\log^2(1/\varepsilon))$ and the cost is
$\mathcal{O}(n\log^2(1/\varepsilon))$.  For the Gaussian
$e^{-\beta y^2/2}$ with $L=1/\sqrt{\beta}$, the pullback
$e^{-\cot^2\theta/2}$ is subgeometric with $\sigma=1/2$, giving
$\mathcal{O}(n\log^2(1/\varepsilon))$.  For the shifted Lorentzian
$1/((y-a)^2+b^2)$ the pullback is meromorphic in $e^{i\theta}$ with
poles at distance $\mathcal{O}(\min(b,L)/L)$ from the real axis,
giving geometric convergence at rate
$d_\delta=\mathcal{O}((L/b)\log(1/\varepsilon))$; centered and at
natural width ($a=0$, $b=L$) this reduces to
$\mathcal{O}(\log(1/\varepsilon))$.

\subsection{Quadrature-weighted preparation}
\label{subsec:quadrature}

Downstream tasks that consume the state as an inner product
$\langle f\mid g\rangle\approx\int_\R f(y)\overline{g(y)}\,dy$ need
the cotangent-grid Jacobian $|dy/d\theta|=L/\sin^2\theta$.  Encoding
the weighted amplitudes $f(y_j)/\sin\theta_j$ builds the Jacobian
in:
\begin{equation}
	\int_{-\infty}^{\infty}\! f(y)\,dy
	=\int_0^\pi\! f(L\cot\theta)\frac{L}{\sin^2\theta}\,d\theta
	\approx\frac{\pi L}{N}\sum_{j=0}^{N-1}\frac{g(\theta_j)}{\sin^2\theta_j}.
	\label{eq:quadrature}
\end{equation}
For the Lorentzian with $L=1$ the weighted amplitude is
$f(y_j)/\sin\theta_j=\sin\theta_j$, again a degree-one trigonometric
polynomial, so the quadrature-weighted preparation also costs
$\mathcal{O}(n)$.  More generally, the weighted variant has the same
asymptotic cost as the unweighted one whenever $g(\theta)/\sin^2\theta$
remains in a class with convergent Fourier expansion, which holds for
$f(y)=\mathcal{O}(y^{-2})$ at infinity.

\subsection{Comparison}
\label{subsec:comparison}

Three families of methods are directly comparable.  The QSVT
pipeline of McArdle, Gily\'en, and
Berta~\cite{McArdleGilyenBerta2025} approximates $f$ on a truncated
interval $[-L_{\mathrm{trunc}},L_{\mathrm{trunc}}]$ through
$y=L_{\mathrm{trunc}}\sin\theta$; the user picks
$L_{\mathrm{trunc}}$ to balance the truncation tail against the
polynomial degree.  The Schmidt-decomposition circuit of
Gonz\'alez-Conde \emph{et al.}~\cite{GonzalezCondeEtAl2020} and the
matrix-product-state construction of Sano and
Hamamura~\cite{SanoHamamura2025} prepare Lorentzians and other
smooth densities at linear depth on a finite domain.
Table~\ref{tab:comparison} collects the gate counts;
Appendix~\ref{app:table} derives each entry under one accounting,
gate count $=n\times$ (polynomial degree) $\times\,1/\mathcal{F}$.

\begin{table*}[t]
	\caption{Gate complexity for preparing a state encoding $f(y)$
		on $n$ qubits to trace distance $\varepsilon$.  The QSVT
		column~\cite{McArdleGilyenBerta2025} is evaluated at the
		truncation radius that balances the tail against the degree;
		for the Lorentzian this is $L_{\mathrm{trunc}}=\varepsilon^{-2/3}$,
		since the truncated-state trace distance scales as the square
		root of the discarded $L^2$-mass.  The linear-depth Lorentzian
		preparations of Refs.~\cite{GonzalezCondeEtAl2020}~(GC)
		and~\cite{SanoHamamura2025}~(SH) are listed where they apply.
		The stereographic method carries no truncation radius; its
		cost is set by the shape-only filling fraction.  Derivations
		are in Appendix~\ref{app:table}.}
	\label{tab:comparison}
	\begin{ruledtabular}
		\begin{tabular}{lllll}
			Target $f(y)$ & QSVT~\cite{McArdleGilyenBerta2025}
			& GC~\cite{GonzalezCondeEtAl2020}
			& SH~\cite{SanoHamamura2025}
			& Stereographic (this work)\\
			\hline
			$1/(1+y^2)$
			& $\mathcal{O}(n\,\varepsilon^{-1}\log\tfrac1\varepsilon)$
			& $\mathcal{O}(n)$ & $\mathcal{O}(n)$ & $\mathcal{O}(n)$\\[2pt]
			$\mathrm{sech}(y)$
			& $\mathcal{O}(n\log^{5/2}\tfrac1\varepsilon)$
			& n.a. & n.a. & $\mathcal{O}(n\log^{2}\tfrac1\varepsilon)$\\[2pt]
			$e^{-\beta y^2/2}$
			& $\tilde{\mathcal{O}}(n\log^{5/4}\tfrac1\varepsilon)$
			& $\mathcal{O}(n)$ & $\mathcal{O}(n)$
			& $\mathcal{O}(n\log^{2}\tfrac1\varepsilon)$\\[2pt]
			$1/((y-a)^2+b^2)$
			& $\mathcal{O}(\tfrac{n}{b}\varepsilon^{-1}\log\tfrac1\varepsilon)$
			& $\mathcal{O}(n)$ & n.a.
			& $\mathcal{O}(\tfrac{nL}{b}\log\tfrac1\varepsilon)$
		\end{tabular}
	\end{ruledtabular}
\end{table*}

For the centered Lorentzian the stereographic method gains a factor
$\varepsilon^{-1}\log(1/\varepsilon)$ over the QSVT pipeline, a
polynomial-in-$\varepsilon$ advantage that comes from the exact
degree-two pullback and the absence of a truncation radius; it ties
the linear-depth preparations of
Refs.~\cite{GonzalezCondeEtAl2020,SanoHamamura2025} at
$\mathcal{O}(n)$, where the stereographic gain is instead
qualitative: one algorithm covers Boyd's class without the
per-target redesign those constructions require, and it removes the
truncation radius.  For the hyperbolic secant it improves on the
QSVT pipeline by $\log^{1/2}(1/\varepsilon)$.  For Gaussians it is
slower than every baseline, by $\log^{3/4}(1/\varepsilon)$ against
the QSVT pipeline and by a logarithmic factor against the
linear-depth methods.  The two-ancilla budget is one to two qubits
below the QSVT pipeline, which uses four in the general case and
three under parity, because the diagonal signal $e^{iA}$ needs no
separate block-encoding ancilla.  Reciprocally, the linear-depth
methods prepare truncated Lorentzians exactly and extend to targets
outside Boyd's class; we claim a complementary route, not a
dominating one.  The output lives on the cotangent grid;
applications requiring a different grid need a regridding step we do
not analyze.  Convergence data for these targets appear in
Fig.~\ref{fig:convergence}.

\begin{figure*}[t]
	\centering
	\includegraphics[width=\textwidth]{figures/convergence_comparison.pdf}
	\caption{$L^\infty$ approximation error versus polynomial degree
		$d$ for the four targets, comparing the rational Chebyshev
		expansion on $\R$ (this work) with truncated Chebyshev
		expansions on $[-L,L]$ at several $L$.  For the Lorentzian
		and shifted Lorentzian the rational basis converges at a
		rate independent of the represented range, while the
		truncated expansions degrade as $L$ grows; for the
		hyperbolic secant and Gaussian the subgeometric rate
		$\exp(-c\sqrt d)$ is visible as the characteristic downward
		curvature, matching the classification of
		Sec.~\ref{subsec:boyd}.}
	\label{fig:convergence}
\end{figure*}

The Lorentzian speedup does not contradict the no-go of
Sec.~\ref{sec:limits}.  That obstruction constrains polynomials of
a block-encoded eigenvalue; here the signal is a walk with
prescribed cotangent-spaced eigenphases and the basis is the
rational Chebyshev family on $\R$.  The chart change alone buys
nothing; the chart paired with the matching grid and signal
operator is what buys the gate savings.
